% !TEX program = xelatex
\documentclass[11pt,a4paper]{article}
\usepackage[margin=2.2cm,top=2cm]{geometry}
\usepackage{fontspec}
\usepackage{polyglossia}
\setdefaultlanguage{turkish}
\usepackage{amsmath,amssymb}
\usepackage{graphicx}
\usepackage{booktabs}
\usepackage{enumitem}
\usepackage{multicol}
\usepackage[most]{tcolorbox}
\usepackage{fancyhdr}
\usepackage{titlesec}
\usepackage{xcolor}
\usepackage{needspace}
\usepackage[hidelinks]{hyperref}

\definecolor{anamavi}{RGB}{0,45,98}
\definecolor{cozumyesil}{RGB}{0,110,50}
\definecolor{uyarikirmizi}{RGB}{180,20,20}

\titleformat{\section}{\large\bfseries\color{anamavi}}{\thesection.}{0.5em}{}
\titleformat{\subsection}{\normalsize\bfseries\color{anamavi}}{\thesubsection}{0.5em}{}

\pagestyle{fancy}\fancyhf{}
\lhead{\small MMAT105 Genel Matematik}
\rhead{\small Hafta 7 -- Ders Notu}
\cfoot{\small\thepage}
\renewcommand{\headrulewidth}{0.4pt}

\newtcolorbox{kural}[1][Kural]{colback=anamavi!5,colframe=anamavi,title=#1,fonttitle=\bfseries,breakable,left=4pt,right=4pt,top=2pt,bottom=2pt}
\newtcolorbox{dikkat}[1][Dikkat]{colback=uyarikirmizi!5,colframe=uyarikirmizi,title=#1,fonttitle=\bfseries,breakable,left=4pt,right=4pt,top=2pt,bottom=2pt}
\newcounter{ornek}
\newenvironment{ornek}{\refstepcounter{ornek}\par\medskip\noindent\textbf{\color{anamavi}Örnek \theornek.}\ }{\par}
\newenvironment{cozum}{\begin{tcolorbox}[colback=cozumyesil!6,colframe=cozumyesil,boxrule=0.5pt,left=4pt,right=4pt,top=2pt,bottom=2pt,breakable]\textbf{\color{cozumyesil}Çözüm.}\ }{\end{tcolorbox}}
\newcounter{alis}
\newenvironment{alistirma}{\refstepcounter{alis}\par\medskip\needspace{9\baselineskip}\noindent\textbf{\color{anamavi}Alıştırma \thealis.}\ }{\par}
\newcommand{\kendinizoncesi}{\needspace{12\baselineskip}}
\newcommand{\ab}[1]{\overline{#1}}
\setlength{\parindent}{0pt}
\setlength{\parskip}{3pt}

% ---------- Sayı doğrusu ----------
\newcommand{\sd}[3]{\begin{tikzpicture}[scale=0.55,baseline=(current bounding box.center)]
\draw[->,thin] (#1-0.7,0)--(#2+0.7,0);
\foreach \t in {#1,...,#2}{\draw (\t,0.12)--(\t,-0.12) node[below=-1pt]{\scriptsize$\t$};}
#3\end{tikzpicture}}
\newcommand{\cizgi}[2]{\draw[anamavi,line width=2.5pt] (#1,0)--(#2,0);}
\newcommand{\isin}[2]{\draw[anamavi,line width=2.5pt,-stealth] (#1,0)--(#2,0);}
\newcommand{\acik}[1]{\draw[anamavi,thick,fill=white] (#1,0) circle (4pt);}
\newcommand{\kapali}[1]{\fill[anamavi] (#1,0) circle (4pt);}

\begin{document}

% ---------------- Başlık ----------------
\begin{minipage}[c]{\textwidth}
{\LARGE\bfseries\color{anamavi} MMAT105 Genel Matematik}\\[4pt]
{\large Hafta 7 Ders Notu: Basit Eşitsizlikler ve Mutlak Değer}\\[6pt]
{\small 2026--2027 Güz}
\end{minipage}

\vspace{4pt}\noindent\textcolor{anamavi}{\rule{\textwidth}{1pt}}

\textbf{Bu notun kullanımı.} Her bölümde önce kısa bir kural kutusu, ardından çözümlü örnekler yer alır. Örnekleri çözümüne bakmadan önce kendiniz deneyin. Notun sonunda önce çözümlü alıştırmalar, ardından cevap anahtarlı ``Kendiniz Çözün'' alıştırmaları yer alır. En sondaki \textbf{karma tekrar} bölümü ara sınava hazırlık içindir; 1--7. haftaların konularını bir arada içerir.

% =====================================================
\section{Eşitsizlik ve Aralık Gösterimi}

\begin{kural}[Semboller ve aralıklar]
\(a<b\) reel sayılar ve \(x\) bir reel sayı olmak üzere:\\[2pt]
\renewcommand{\arraystretch}{1.25}
\begin{tabular}{@{}lll@{\qquad}lll@{}}
\(a<x<b\) & \((a,b)\) & açık aralık & \(x\ge a\) & \([a,\infty)\) & \\
\(a\le x\le b\) & \([a,b]\) & kapalı aralık & \(x>a\) & \((a,\infty)\) & \\
\(a\le x<b\) & \([a,b)\) & yarı açık & \(x\le a\) & \((-\infty,a]\) & \\
\(a<x\le b\) & \((a,b]\) & yarı açık & \(x<a\) & \((-\infty,a)\) &
\end{tabular}\\[2pt]
Sayı doğrusunda dahil olan uç dolu yuvarlak (\(\bullet\)), dahil olmayan uç boş yuvarlak (\(\circ\)) ile gösterilir.
\end{kural}

\begin{dikkat}
\(\infty\) bir sayı değildir; \(-\infty\) ve \(\infty\) yanında parantez daima açıktır. \ \(\mathbb{R}=(-\infty,\infty)\).
\end{dikkat}

\begin{ornek}
\(A=[-2,3)\) ve \(B=(0,5]\) aralıkları veriliyor. \(A\) ile \(B\)'nin ortak kısmını (kesişimini) ve birleşimini aralık olarak yazınız.
\end{ornek}
\begin{cozum}
\(A\): \sd{-3}{6}{\cizgi{-2}{3}\kapali{-2}\acik{3}} \qquad \(B\): \sd{-3}{6}{\cizgi{0}{5}\acik{0}\kapali{5}}\\[3pt]
Ortak kısım: hem \(x<3\) hem \(x>0\) \(\Rightarrow\) \(A\cap B=(0,3)\). \quad Birleşim: \(A\cup B=[-2,5]\).
\end{cozum}

% =====================================================
\newpage
\section{Eşitsizliğin Özellikleri ve Birinci Dereceden Eşitsizlikler}

\begin{kural}[Özellikler]
\(a<b\) olsun.
\begin{itemize}[leftmargin=*,itemsep=1pt]
\item Her iki tarafa aynı sayı eklenir veya çıkarılır: \(a+c<b+c\).
\item \(c>0\) ile çarpma/bölmede yön korunur: \(ac<bc\). \quad \(c<0\) ile çarpma/bölmede yön \textbf{değişir}: \(ac>bc\).
\item Aynı yönlü eşitsizlikler taraf tarafa \textbf{toplanır}: \(a<b\) ve \(c<d\Rightarrow a+c<b+d\). Taraf tarafa çıkarma veya bölme yapılamaz.
\item \(a\) ile \(b\) aynı işaretli ise (\(0<a<b\) veya \(a<b<0\)): \(\dfrac1a>\dfrac1b\).
\item \(ax+b<0\) tipindeki eşitsizlik denklem gibi çözülür; negatif sayıyla çarpma/bölmede yön çevrilir. Çift eşitsizlikte (\(m<ax+b\le n\)) işlemler üç tarafa birden uygulanır.
\end{itemize}
\end{kural}

\begin{dikkat}[Sık yapılan hatalar]
\(-3x<12\Rightarrow x>-4\) (yön değişir). \quad \(5<6\) ve \(1<4\) iken \(5-1<6-4\) \emph{yanlıştır}. \quad \(x-3<x+1\) gibi \(x\)'lerin sadeleştiği durumda sonuç ya \(\mathbb{R}\) ya da \(\varnothing\) olur.
\end{dikkat}

\begin{ornek}
\(a<b<0\) olduğuna göre \(a^2\) ile \(b^2\) ve \(\dfrac1a\) ile \(\dfrac1b\) sayılarını sıralayınız.
\end{ornek}
\begin{cozum}
\(a<b<0\) ise \(a\), sıfıra \(b\)'den daha uzaktır; karesi daha büyüktür: \(a^2>b^2\).\\
\(a\) ve \(b\) aynı işaretli olduğundan ters alınınca yön döner: \(\dfrac1a>\dfrac1b\).\\
Kontrol: \(a=-3,\ b=-2\): \(9>4\) ve \(-\tfrac13>-\tfrac12\) \checkmark
\end{cozum}

\begin{ornek}
\(\dfrac{2x+1}{3}-\dfrac{x-2}{4}\ge 2\) eşitsizliğini sağlayan en küçük tam sayı kaçtır?
\end{ornek}
\begin{cozum}
Her terimi \(12\) ile çarpalım: \(4(2x+1)-3(x-2)\ge 24\Rightarrow 8x+4-3x+6\ge 24\Rightarrow 5x\ge 14\Rightarrow x\ge \dfrac{14}{5}=2{,}8\).\\
En küçük tam sayı: \(3\).
\end{cozum}

\begin{ornek}
\(2(x-3)-5>4x+1\) eşitsizliğinin çözüm kümesi nedir? Bu eşitsizliği sağlayan en büyük tam sayı kaçtır?
\end{ornek}
\begin{cozum}
\(2x-11>4x+1\Rightarrow -2x>12\). \(-2\)'ye bölünce yön değişir: \(x<-6\).\\
\(\text{ÇK}=(-\infty,-6)\); \ en büyük tam sayı \(-7\).
\end{cozum}

\begin{ornek}
\(-4<5-3x\le 11\) eşitsizliğinin çözüm kümesini bulunuz ve sayı doğrusunda gösteriniz.
\end{ornek}
\begin{cozum}
Üç taraftan \(5\) çıkaralım: \(-9<-3x\le 6\). \(-3\)'e bölünce iki yön de değişir: \(3>x\ge -2\), yani \(-2\le x<3\).\\
\(\text{ÇK}=[-2,3)\) \qquad \sd{-3}{4}{\cizgi{-2}{3}\kapali{-2}\acik{3}}
\end{cozum}

\begin{ornek}
\(3x-1>5\) ve \(2x+3\le 15\) eşitsizliklerini \textbf{birlikte} sağlayan kaç tam sayı vardır?
\end{ornek}
\begin{cozum}
\(3x>6\Rightarrow x>2\); \quad \(2x\le 12\Rightarrow x\le 6\). Ortak çözüm: \(2<x\le 6\), yani \((2,6]\).\\
Tam sayılar: \(3,4,5,6\) \(\Rightarrow\) \(4\) tane.
\end{cozum}

% =====================================================
\section{Tam Sayı Değerleri ve Değer Aralığı}

\begin{kural}[Sayma ve aralık kurma]
\begin{itemize}[leftmargin=*,itemsep=1pt]
\item \(a<b\) tam sayıları için: \([a,b]\)'de \(b-a+1\), \ \((a,b)\)'de \(b-a-1\), \ \([a,b)\) ve \((a,b]\)'de \(b-a\) tam sayı vardır.
\item Bir ifadenin aralığı, \(x\)'in aralığından adım adım kurulur (negatifle çarpınca yön değişir).
\item \(x+y\): aralıklar toplanır. \ \(x-y\): \(-y\)'nin aralığı bulunur, \(x\)'inkiyle toplanır.
\item \(x\cdot y\): uç değerlerin dört çarpımından en küçüğü ve en büyüğü alınır.
\item \(x^2\): aralık \(0\)'ı içeriyorsa en küçük değer \(0\)'dır (dahil); en büyük değer, sıfıra en uzak ucun karesidir.
\end{itemize}
\end{kural}

\begin{dikkat}
\(x,y\) \textbf{tam sayı} ise önce alabilecekleri değerler listelenir; en büyük/en küçük değer bu listelerden seçilir. \(x\) reel ise uç dahil değilse o değer alınamaz.
\end{dikkat}

\begin{ornek}
\(-1\le x<3\) ise \(4-2x\) ifadesi hangi aralıktadır? Kaç farklı tam sayı değeri alabilir?
\end{ornek}
\begin{cozum}
\(-2\) ile çarpalım: \(2\ge -2x>-6\), yani \(-6<-2x\le 2\). \(4\) ekleyelim: \(-2<4-2x\le 6\), yani \((-2,6]\).\\
Tam sayı değerleri: \(-1,0,\dots,6\) \(\Rightarrow\) \(6-(-2)=8\) tane.
\end{cozum}

\begin{ornek}
\(1<x<4\) ve \(2<y<5\) ise \(2x-3y\) hangi aralıktadır? \(x\) ve \(y\) tam sayı ise \(2x-3y\)'nin en büyük ve en küçük değeri kaçtır?
\end{ornek}
\begin{cozum}
\(2<2x<8\) ve \(-15<-3y<-6\); toplarsak \(-13<2x-3y<2\).\\
Tam sayılar için \(x\in\{2,3\}\), \(y\in\{3,4\}\): en büyük \(2\cdot3-3\cdot3=-3\), en küçük \(2\cdot2-3\cdot4=-8\).
\end{cozum}

\begin{ornek}
\(-2<x<3\) ve \(1<y<4\) ise \(x\cdot y\) hangi aralıktadır?
\end{ornek}
\begin{cozum}
Uç çarpımlar: \((-2)\cdot1=-2\), \ \((-2)\cdot4=-8\), \ \(3\cdot1=3\), \ \(3\cdot4=12\). \ En küçük \(-8\), en büyük \(12\); uçlar dahil değil:
\[-8<x\cdot y<12\]
\end{cozum}

% =====================================================
\section{Mutlak Değer: Tanım ve Özellikler}

\begin{kural}[Tanım ve özellikler]
\[|x|=\begin{cases}x, & x\ge 0\\ -x, & x<0\end{cases}\qquad |x-a|:\ x \text{ ile } a \text{ arasındaki uzaklık}\]
\(|x|\ge 0\), \ \(|-x|=|x|\), \ \(|x-y|=|y-x|\), \ \(|x\cdot y|=|x|\cdot|y|\), \ \(\left|\dfrac xy\right|=\dfrac{|x|}{|y|}\ (y\neq0)\). \ Üçgen eşitsizliği: \(|x+y|\le|x|+|y|\).\\
\textbf{İşareti bilinen ifade:} \(A\ge 0\) ise \(|A|=A\); \ \(A<0\) ise \(|A|=-A\).
\end{kural}

\begin{dikkat}
\(|-x|=x\) yalnızca \(x\ge0\) iken doğrudur. \quad \(|3-x|\) içerisi negatifse \(-(3-x)=x-3\) yazılır; parantezi unutmayın.
\end{dikkat}

\needspace{6\baselineskip}
\begin{ornek}
\(a=-3\) ve \(b=5\) için \(|a|-|b|+|a-b|-|b-a|\) kaçtır?
\end{ornek}
\begin{cozum}
\(|a-b|=|b-a|\) olduğundan son iki terim birbirini götürür: \(|-3|-|5|=3-5=-2\).
\end{cozum}

\begin{ornek}
\(a<b<0<c\) olduğuna göre \(|a-b|+|c-a|-|b-c|\) ifadesinin eşiti nedir?
\end{ornek}
\begin{cozum}
\(a-b<0\Rightarrow|a-b|=b-a\); \quad \(c-a>0\Rightarrow|c-a|=c-a\); \quad \(b-c<0\Rightarrow|b-c|=c-b\).
\[(b-a)+(c-a)-(c-b)=2b-2a\]
Kontrol: \(a=-3,\ b=-1,\ c=2\): \(2+5-3=4\) ve \(2(-1)-2(-3)=4\) \checkmark
\end{cozum}

\begin{ornek}
\(2<x<5\) olduğuna göre \(|x-2|+|x-5|-|x|\) ifadesinin eşiti nedir?
\end{ornek}
\begin{cozum}
\(x-2>0\), \(x-5<0\), \(x>0\): \quad \((x-2)+(5-x)-x=3-x\). \quad Kontrol: \(x=3\): \(1+2-3=0=3-3\) \checkmark
\end{cozum}

% =====================================================
\section{Mutlak Değerli Denklemler}

\begin{kural}[Denklemler]
\(|ax+b|=c\): \ \(c>0\) ise \(ax+b=c\) veya \(ax+b=-c\); \ \(c=0\) ise \(ax+b=0\); \ \(c<0\) ise \(\text{ÇK}=\varnothing\).\\
\(|f(x)|=|g(x)|\ \Rightarrow\ f(x)=g(x)\) veya \(f(x)=-g(x)\).\\
Sağ taraf \(x\) içeriyorsa (\(|f(x)|=g(x)\)) bulunan kökler denklemde \textbf{mutlaka kontrol edilir}.
\end{kural}

\begin{ornek}
\(|3x-4|=5\) denkleminin kökleri çarpımı kaçtır?
\end{ornek}
\begin{cozum}
\(3x-4=5\Rightarrow x=3\); \quad \(3x-4=-5\Rightarrow x=-\dfrac13\). \quad Kökler çarpımı: \(3\cdot\left(-\dfrac13\right)=-1\).
\end{cozum}

\begin{ornek}
\(|2x-1|=|x+4|\) denkleminin kökleri toplamı kaçtır?
\end{ornek}
\begin{cozum}
\(2x-1=x+4\Rightarrow x=5\); \quad \(2x-1=-(x+4)\Rightarrow 3x=-3\Rightarrow x=-1\).\\
Kontrol: \(|9|=|9|\), \ \(|-3|=|3|\). \ Kökler toplamı: \(5+(-1)=4\).
\end{cozum}

\begin{ornek}
\(|x-2|+2x=7\) denkleminin çözüm kümesi nedir?
\end{ornek}
\begin{cozum}
\(x\ge2\) ise \(|x-2|=x-2\): \ \(x-2+2x=7\Rightarrow x=3\) \ (\(3\ge2\), uygun).\\
\(x<2\) ise \(|x-2|=2-x\): \ \(2-x+2x=7\Rightarrow x=5\) \ (\(5<2\) değil, \textbf{uygun değil}).\\
\(\text{ÇK}=\{3\}\). \ Kontrol: \(|1|+6=7\) \checkmark
\end{cozum}

% =====================================================
\section{Mutlak Değerli Eşitsizlikler}

\begin{kural}[Eşitsizlikler]
\(b>0\) olmak üzere:\\
\(|x-a|<b\ \Rightarrow\ a-b<x<a+b\) \qquad (\(a\)'ya uzaklığı \(b\)'den küçük olanlar)\\
\(|x-a|>b\ \Rightarrow\ x<a-b\) \ veya \ \(x>a+b\) \qquad (\(a\)'ya uzaklığı \(b\)'den büyük olanlar)\\
\(|ax+c|\le b\ \Rightarrow\ -b\le ax+c\le b\). \quad \(m\le|f(x)|<n\) tipinde iki eşitsizlik ayrı çözülüp ortak kısım alınır.
\end{kural}

\begin{dikkat}
\(|f(x)|<\) negatif sayı \(\Rightarrow\varnothing\); \quad \(|f(x)|\ge\) negatif sayı \(\Rightarrow\mathbb{R}\); \quad \(|f(x)|\le 0\Rightarrow f(x)=0\).
\end{dikkat}

\begin{ornek}
\(|x-4|\le 3\) eşitsizliğini sağlayan tam sayıların sayısı ve toplamı kaçtır?
\end{ornek}
\begin{cozum}
\(-3\le x-4\le 3\Rightarrow 1\le x\le 7\). \ Tam sayılar \(1,\dots,7\): \(7\) tane; \ toplam \(\dfrac{7\cdot 8}{2}=28\).
\end{cozum}

\begin{ornek}
\(-10\le x\le 10\) olan tam sayılardan kaç tanesi \(|2x+3|>5\) eşitsizliğini \textbf{sağlamaz}?
\end{ornek}
\begin{cozum}
Sağlamayanlar \(|2x+3|\le 5\) olanlardır: \(-5\le 2x+3\le 5\Rightarrow -8\le 2x\le 2\Rightarrow -4\le x\le 1\).\\
Tam sayılar \(-4,-3,-2,-1,0,1\) \(\Rightarrow\) \(6\) tane.
\end{cozum}

\begin{ornek}
\(1<|x+2|\le 4\) eşitsizliğini sağlayan tam sayıların toplamı kaçtır?
\end{ornek}
\begin{cozum}
\(|x+2|\le 4\Rightarrow -6\le x\le 2\). \quad \(|x+2|>1\Rightarrow x<-3\) veya \(x>-1\).\\
Ortak tam sayılar: \(-6,-5,-4,\ 0,1,2\). \ Toplam: \(-15+3=-12\).
\end{cozum}

\begin{ornek}
\(|x-a|<b\) eşitsizliğinin çözüm kümesi \((-3,9)\) ise \(a\) ve \(b\) kaçtır?
\end{ornek}
\begin{cozum}
Çözüm \(a-b<x<a+b\) biçimindedir: \(a-b=-3\), \(a+b=9\). Taraf tarafa toplarsak \(2a=6\Rightarrow a=3\), \(b=6\).\\
Yorum: \(a\), aralığın orta noktası; \(b\), aralık uzunluğunun yarısıdır.
\end{cozum}

% =====================================================
\section{Özet: Formül Kartı}
\begin{center}
\renewcommand{\arraystretch}{1.35}
\begin{tabular}{@{}ll@{}}
\toprule
Konu & Formül / Kural\\\midrule
Aralık & \(a<x\le b\Leftrightarrow (a,b]\); \ \(\infty\) yanında parantez açık\\
Negatifle çarpma/bölme & eşitsizliğin yönü değişir\\
Taraf tarafa & aynı yönlüler toplanır; çıkarılmaz, bölünmez\\
Ters alma & \(a,b\) aynı işaretli, \(a<b\Rightarrow \frac1a>\frac1b\)\\
Tam sayı sayısı & \([a,b]\): \(b-a+1\); \ \((a,b)\): \(b-a-1\); \ yarı açık: \(b-a\)\\
İki değişken & \(x-y\): \(x+(-y)\); \ \(x\cdot y\): dört uç çarpımı; \ \(0\) aralıkta ise \(\min x^2=0\)\\
Mutlak değer & \(|A|=A\ (A\ge0)\), \ \(|A|=-A\ (A<0)\); \ \(|x-a|\): uzaklık\\
\(|ax+b|=c\) & \(c>0\): iki kök; \ \(c=0\): tek kök; \ \(c<0\): \(\varnothing\)\\
\(|x-a|<b\) & \(a-b<x<a+b\)\\
\(|x-a|>b\) & \(x<a-b\) veya \(x>a+b\)\\\bottomrule
\end{tabular}
\end{center}

% =====================================================
\needspace{15\baselineskip}
% Hafta 7 -- alıştırma gövdesi (ders notu ve ayrı alıştırma kağıdı tarafından \input edilir)
% Sayı doğrusu makroları (\sd, \cizgi, \isin, \acik, \kapali) her iki belgenin preamble'ında tanımlıdır.
\section{Çözümlü Alıştırmalar}
Her soruyu önce kendiniz deneyin, sonra çözümü okuyun.

\begin{alistirma}
\(\dfrac{x+1}{2}-\dfrac{x-3}{5}<2\) eşitsizliğini sağlayan en büyük tam sayı kaçtır?
\end{alistirma}
\begin{cozum}
Paydaların EKOK'u \(10\); her terimi \(10\) ile çarpalım (\(10>0\), yön korunur):
\[5(x+1)-2(x-3)<20\ \Rightarrow\ 5x+5-2x+6<20\ \Rightarrow\ 3x+11<20\ \Rightarrow\ 3x<9\ \Rightarrow\ x<3\]
\(3\) dahil değil \(\Rightarrow\) en büyük tam sayı \(2\).
\end{cozum}

\begin{alistirma}
\(4-x\ge 2x-11\) ve \(x>-3\) eşitsizliklerini birlikte sağlayan kaç tam sayı vardır?
\end{alistirma}
\begin{cozum}
\(4+11\ge 2x+x\Rightarrow 15\ge 3x\Rightarrow x\le 5\). \ İkinci koşul \(x>-3\). Ortak çözüm: \(-3<x\le 5\), yani \((-3,5]\).\\
\sd{-4}{6}{\cizgi{-3}{5}\acik{-3}\kapali{5}}\\
Tam sayılar: \(-2,-1,0,1,2,3,4,5\) \(\Rightarrow\) \(5-(-3)=8\) tane.
\end{cozum}

\begin{alistirma}
\(3<x<7\) ve \(-2<y<4\) ise \(x+y\) ve \(x-y\) hangi aralıklardadır?
\end{alistirma}
\begin{cozum}
\(x+y\): aralıklar taraf tarafa toplanır: \(3+(-2)<x+y<7+4\Rightarrow 1<x+y<11\).\\
\(x-y\): önce \(y\)'nin aralığını \(-1\) ile çarpalım (yön değişir): \(-4<-y<2\). Sonra \(x\) ile toplayalım:
\[3+(-4)<x-y<7+2\ \Rightarrow\ -1<x-y<9\]
\end{cozum}

\begin{alistirma}
\(-3<x<2\), \ \(-1<y<5\) ve \(x,y\) tam sayılardır. \(x\cdot y\)'nin en küçük ve en büyük değeri kaçtır?
\end{alistirma}
\begin{cozum}
\(x\in\{-2,-1,0,1\}\), \quad \(y\in\{0,1,2,3,4\}\) (\(y\) hiç negatif değil).\\
En küçük: çarpımın negatif olması için \(x\) negatif olmalı; \(|x|\) ve \(y\) en büyük seçilir: \(x=-2\), \(y=4\) \(\Rightarrow\) \(-8\).\\
En büyük: çarpım pozitif olmalı; bunun için \(x=1\) (tek pozitif değer) ve \(y=4\): \(1\cdot 4=4\).
\end{cozum}

\begin{alistirma}
\(a<0<b\) olduğuna göre \(|a-b|-|b|-|2a|\) ifadesinin eşiti nedir?
\end{alistirma}
\begin{cozum}
\(a-b<0\Rightarrow |a-b|=b-a\); \quad \(b>0\Rightarrow |b|=b\); \quad \(2a<0\Rightarrow |2a|=-2a\).
\[(b-a)-b-(-2a)=b-a-b+2a=a\]
Kontrol: \(a=-1,\ b=2\): \(|-3|-|2|-|-2|=3-2-2=-1=a\) \checkmark
\end{cozum}

\begin{alistirma}
\(|x+2|=|x-6|\) denkleminin çözüm kümesi nedir?
\end{alistirma}
\begin{cozum}
\(x+2=x-6\Rightarrow 2=-6\): olanaksız, bu durumdan kök gelmez.\\
\(x+2=-(x-6)\Rightarrow x+2=-x+6\Rightarrow 2x=4\Rightarrow x=2\). \quad Kontrol: \(|4|=|-4|\) \checkmark \quad \(\text{ÇK}=\{2\}\).\\
Uzaklık yorumu: \(x\)'in \(-2\)'ye ve \(6\)'ya uzaklıkları eşit \(\Rightarrow\) \(x\) orta noktadır: \(\dfrac{-2+6}{2}=2\).
\end{cozum}

\begin{alistirma}
\(-5\le x\le 5\) olan tam sayılardan kaç tanesi \(|3x-1|\ge 8\) eşitsizliğini sağlar?
\end{alistirma}
\begin{cozum}
\(3x-1\ge 8\Rightarrow 3x\ge 9\Rightarrow x\ge 3\) \quad veya \quad \(3x-1\le -8\Rightarrow 3x\le -7\Rightarrow x\le -\dfrac73\approx-2{,}33\).\\
\([-5,5]\)'teki tam sayılardan: \(x\ge 3\) için \(3,4,5\); \ \(x\le -\frac73\) için \(-5,-4,-3\) (\(-2\) sağlamaz: \(|-7|=7<8\)).\\
Toplam \(6\) tane.
\end{cozum}

\begin{alistirma}
\(2\le|x-3|\le 5\) eşitsizliğini sağlayan tam sayıların sayısı ve toplamı kaçtır?
\end{alistirma}
\begin{cozum}
\(|x-3|\le 5\Rightarrow -5\le x-3\le 5\Rightarrow -2\le x\le 8\).\\
\(|x-3|\ge 2\Rightarrow x-3\ge 2\) veya \(x-3\le -2\) \(\Rightarrow\) \(x\ge 5\) veya \(x\le 1\).\\
Ortak kısım: \([-2,1]\cup[5,8]\) \qquad \sd{-3}{9}{\cizgi{-2}{1}\cizgi{5}{8}\kapali{-2}\kapali{1}\kapali{5}\kapali{8}}\\
Tam sayılar: \(-2,-1,0,1,5,6,7,8\) \(\Rightarrow\) \(8\) tane; \ toplam \((-2)+26=24\).
\end{cozum}

\kendinizoncesi
\section{Alıştırmalar (Kendiniz Çözün)}

\begin{enumerate}[label=\textbf{\arabic*.},leftmargin=*,itemsep=4pt]
\item \(\{x\in\mathbb{R}\,:\,-3\le x<5\}\) kümesini aralık gösterimiyle yazınız. \((-\infty,2]\) aralığını eşitsizlikle ifade ediniz.
\item \(5x-3\le 2x+9\) eşitsizliğinin çözüm kümesi nedir?
\item \(7-2x>3(x-1)\) eşitsizliğinin çözüm kümesi nedir?
\item \(-2x+7\ge 1\) eşitsizliğini sağlayan kaç tane doğal sayı vardır? \ (\(\mathbb{N}=\{0,1,2,\dots\}\))
\item \(-5<2x+3\le 11\) eşitsizliğini sağlayan kaç tam sayı vardır? Bunların toplamı kaçtır?
\item \(-1<a<3\) ise \(5-2a\) hangi aralıktadır?
\item \(-4<x<3\) ise \(x^2\) hangi aralıktadır?
\item \(|-6|-|4-9|+|-2|\cdot|3|\) işleminin sonucu kaçtır?
\item \(x<-1\) olduğuna göre \(|x+1|+|x-3|\) ifadesinin eşiti nedir?
\item \(|4x-2|=10\) denkleminin çözüm kümesi nedir?
\item \(|x-5|<3\) eşitsizliğini sağlayan tam sayıların sayısı ve toplamı kaçtır?
\item \(|x-a|\le b\) eşitsizliğinin çözüm kümesi \([-1,7]\) ise \(a\cdot b\) kaçtır?
\end{enumerate}

\subsection*{Cevap Anahtarı}
{\small
\begin{multicols}{2}
\begin{enumerate}[label=\textbf{\arabic*.},leftmargin=*,itemsep=1pt]
\item \([-3,5)\); \quad \(x\le 2\).
\item \(3x\le 12\Rightarrow x\le 4\); \ \((-\infty,4]\).
\item \(10>5x\Rightarrow x<2\); \ \((-\infty,2)\).
\item \(-2x\ge -6\Rightarrow x\le 3\): \(0,1,2,3\) \(\Rightarrow\) \(4\) tane.
\item \(-4<x\le 4\): \(-3,\dots,4\); \(8\) tane, toplam \(4\).
\item \(-6<-2a<2\Rightarrow -1<5-2a<7\); \ \((-1,7)\).
\item \(0\) aralıkta, en uzak uç \(-4\): \([0,16)\).
\item \(6-5+2\cdot 3=7\).
\item \(-(x+1)+(3-x)=2-2x\).
\item \(4x-2=\pm10\Rightarrow x=3\) veya \(x=-2\); \ \(\{-2,3\}\).
\item \(2<x<8\): \(3,\dots,7\); \(5\) tane, toplam \(25\).
\item \(a-b=-1\), \(a+b=7\Rightarrow a=3,\ b=4\); \ \(a\cdot b=12\).
\end{enumerate}
\end{multicols}}

% =====================================================
\newcounter{karma}
\newenvironment{karma}[1]{\refstepcounter{karma}\par\vspace{2pt}\needspace{7\baselineskip}\noindent\textbf{\color{anamavi}K\thekarma.}\ \textit{(Hafta #1)}\ }{\par}

\needspace{15\baselineskip}
\section{Ara Sınava Hazırlık: Karma Tekrar (Hafta 1--7)}
K5, K6, K8 ve K12'nin çözümü sorunun altında, diğerlerinin cevabı bölüm sonundadır.

\begin{karma}{1}
\(5-2\cdot(3-7)+(-2)^3\div 4\) işleminin sonucu kaçtır?
\end{karma}

\begin{karma}{1}
\(\tfrac{8!}{6!\cdot 2!}+0!\) işleminin sonucu kaçtır?
\end{karma}

\begin{karma}{2}
\(360\) sayısının kaç tane pozitif tam böleni vardır?
\end{karma}

\begin{karma}{2}
Bir doğal sayı \(7\)'ye bölündüğünde bölüm \(12\)'dir. Kalan alabileceği en büyük değeri aldığında bölünen sayı kaçtır?
\end{karma}

\begin{karma}{3}
Dört basamaklı \(\ab{3a5b}\) sayısı \(45\) ile tam bölünmektedir. \(a\)'nın alabileceği değerlerin toplamı kaçtır?
\end{karma}
\begin{cozum}
\(45=5\cdot 9\): \(5\)'e bölünme \(\Rightarrow b\in\{0,5\}\); \ \(9\)'a bölünme \(\Rightarrow 8+a+b\) dokuzun katı.\\
\(b=0\): \(8+a\in\{9,18\}\Rightarrow a=1\) (\(a=10\) rakam değil); \ \(b=5\): \(13+a=18\Rightarrow a=5\) \ (\(3150\), \(3555\)).\\
\(a\) değerlerinin toplamı: \(1+5=6\).
\end{cozum}

\begin{karma}{3}
\(18\) ile de \(24\) ile de bölündüğünde \(5\) kalanını veren en küçük iki basamaklı doğal sayı kaçtır?
\end{karma}
\begin{cozum}
\(N-5\), \(18=2\cdot 3^2\) ve \(24=2^3\cdot 3\)'ün ortak katı: \(\text{EKOK}=2^3\cdot 3^2=72\), \ \(N=72k+5\).\\
\(k=0\Rightarrow N=5\) (tek basamaklı); \ \(k=1\Rightarrow N=77\). \quad Kontrol: \(77=18\cdot 4+5=24\cdot 3+5\) \checkmark
\end{cozum}

\begin{karma}{4}
\(\dfrac{0{,}6+\frac25}{1{,}25-\frac34}\) işleminin sonucu kaçtır?
\end{karma}

\begin{karma}{4}
\(a=\dfrac34\), \(b=\dfrac57\), \(c=\dfrac79\) sayılarını küçükten büyüğe sıralayınız.
\end{karma}
\begin{cozum}
\(1\)'e uzaklıklar: \(1-a=\frac14=\frac28\), \ \(1-b=\frac27\), \ \(1-c=\frac29\). Paylar eşitken paydası büyük olan küçüktür: \(\frac29<\frac28<\frac27\).\\
\(1\)'e en yakın olan en büyüktür \(\Rightarrow\) \(b<a<c\). \ (Kontrol: \(a=\frac{189}{252}\), \(b=\frac{180}{252}\), \(c=\frac{196}{252}\).)
\end{cozum}

\begin{karma}{5}
\(\dfrac{x-2}{3}+\dfrac{x+1}{2}=4\) denkleminin çözümü nedir?
\end{karma}

\begin{karma}{6}
\(2x+y=11\) ve \(x-y=1\) ise \(x\cdot y\) kaçtır?
\end{karma}

\begin{karma}{7}
\(-3<2x-5\le 7\) eşitsizliğini sağlayan kaç tam sayı vardır?
\end{karma}

\begin{karma}{7}
\(|2x-3|\le 5\) eşitsizliğini sağlayan tam sayıların toplamı kaçtır?
\end{karma}
\begin{cozum}
\(-5\le 2x-3\le 5\). Üç tarafa \(3\) ekleyelim: \(-2\le 2x\le 8\); \(2\)'ye bölelim: \(-1\le x\le 4\).\\
Tam sayılar: \(-1,0,1,2,3,4\). \quad Toplam: \(-1+0+1+2+3+4=9\).
\end{cozum}

\subsection*{Karma Tekrar Cevapları (çözümü verilmeyenler)}
{\small
\begin{multicols}{2}
\textbf{K1.} \(5-2\cdot(-4)+(-8)\div 4=5+8-2=11\).\par
\textbf{K2.} \(\tfrac{8\cdot 7}{2}+1=28+1=29\).\par
\textbf{K3.} \(360=2^3\cdot 3^2\cdot 5\): \((3+1)(2+1)(1+1)=24\).\par
\textbf{K4.} En büyük kalan \(6\): \(7\cdot 12+6=90\).\par
\textbf{K7.} \((0{,}6+0{,}4)\div(1{,}25-0{,}75)=1\div 0{,}5=2\).\par
\textbf{K9.} \(6\) ile çarp: \(2x-4+3x+3=24\Rightarrow x=5\).\par
\textbf{K10.} Topla: \(3x=12\Rightarrow x=4\), \(y=3\); \(x\cdot y=12\).\par
\textbf{K11.} \(2<2x\le 12\Rightarrow 1<x\le 6\): \(5\) tane.\par
\end{multicols}}


\vspace{6pt}
\noindent\textcolor{anamavi}{\rule{\textwidth}{0.5pt}}\\
{\small Gelecek hafta: ARA SINAV (Hafta 1--7 konuları).}

\end{document}
